%!TEX root = ../thesis.tex

\section{Context within thesis}

The previous chapter introduced the parallel pathway theory (PPT) as a circuit mechanism for systems memory consolidation. In that model, parallel pathways learn with different effective learning rates, leading to different degrees of semantic generalization. The present chapter asks what such a diversity of learning rates could be useful for when the statistics of the environment change over time. The chapter uses Bayesian change-point detection as a normative framework for learning in changing environments, and suggests that a diversity in learning rates during systems consolidation could be part of a neural implementation thereof.

\section{Introduction}

\todo[inline]{better citation for McClelland1995}

Memories allow us to predict the future based on the past. This requires the integration of previous experiences into an internal model of the world. Such integration is thought to occur during systems memory consolidation, where new experiences stored in the hippocampus are gradually integrated into neocortical knowledge structures \citep{McClelland1995}. This is especially relevant when learning the statistical regularities of the environment. However, the world and its statistics are not stationary. When the environment changes, old experiences can become misleading, and predictions that were useful before may no longer be appropriate. For example, if the commute to work used to be reliably quick but then roadworks start, the previous distribution of travel times is no longer the right basis for predicting how long the commute will take tomorrow.

A principled solution to this problem is provided by Bayesian change-point detection algorithms \citep{Adams2007}. These algorithms infer whether a change in the statistics of the environment has occurred and adjust predictions by weighting evidence from different possible intervals since the last change. This allows an agent to combine two requirements that are in tension: integrating evidence over long periods when the world is stable, but rapidly discounting old evidence when the world changes.

Exact Bayesian change-point detection is computationally expensive because it requires maintaining predictive distributions for different possible run lengths. Several approximations have therefore been proposed. \citet{Nassar2010} reduced Bayesian updating to a delta-rule update with an adaptive learning rate. A mixture of fixed-learning-rate estimators can also approximate the full Bayesian solution \citep{Wilson2013}. This latter approximation is especially relevant here, because systems memory consolidation theory suggests different learning rates in the consolidation hierarchy \citep{McClelland1995,Roxin2013}, and PPT introduces a learning-rate dependent semantic generalization strength (Ch.~\ref{ch:ppt}). In this chapter, we ask whether learning-rate diversity in systems consolidation could provide a biological basis for the multiple evidence-window estimates required by approximate change-point inference.

\section{Learning a Bernoulli contingency in a changing environment}

We consider the simplest non-trivial case: a binary observation $x_t \in \{0,1\}$ at trial $t$, drawn from a Bernoulli distribution with success probability $\theta_t$,
\begin{equation}
    p(x_t \mid \theta_t) = \theta_t^{x_t}(1-\theta_t)^{1-x_t}.
\end{equation}
The latent parameter $\theta_t$ is piecewise constant: it takes a fixed value during each stable regime and resets at change points (Fig.~\ref{fig:cpd-bernoulli-task}A). The learner only observes the binary sequence $x_{1:t} = (x_1, \dots, x_t)$ (Fig.~\ref{fig:cpd-bernoulli-task}B) and has to predict the next observation $x_{t+1}$ without knowing the change points.

If the environment were stationary, $\theta_t$ would be constant and a Beta prior with pseudo-counts $a_0$ and $b_0$ would give a conjugate Beta posterior,
\begin{equation}
    p(\theta) = \frac{1}{B(a_0,b_0)}\theta^{a_0-1}(1-\theta)^{b_0-1},
\end{equation}
\begin{equation}
    p(\theta \mid x_{1:t})
    = \frac{1}{B(a_t,b_t)}\theta^{a_t-1}(1-\theta)^{b_t-1},
\end{equation}
where $B(\cdot,\cdot)$ is the beta function and the running counts
\begin{equation}
    a_t = a_0 + \sum_{i=1}^{t}x_i,
    \qquad
    b_t = b_0 + \sum_{i=1}^{t}(1-x_i)
\end{equation}
accumulate successes and failures over the full history. The posterior predictive mean is then $\hat\theta_t = a_t/(a_t+b_t)$. This stationary learner is useful when the environment is genuinely stable, but it averages over all past observations and therefore adapts too slowly when the contingency changes (Fig.~\ref{fig:cpd-bernoulli-task}C, red line).

\begin{figure}[htbp]
\centering
\includegraphics[width=0.92\linewidth]{cpd_bernoulli_task.pdf}
\caption{Learning a Bernoulli contingency in a changing environment. Panel A: latent Bernoulli parameter $\theta_t$ over trials, with vertical lines marking change points. Panel B: observed binary sequence $x_t$. Panel C: stationary Beta learner accumulates all evidence and therefore adapts slowly after changes, while BOCPD adapts after inferred change points.}
\label{fig:cpd-bernoulli-task}
\end{figure}

\section{Bayesian online change-point detection}
\label{sec:cpd-bocpd}

Bayesian online change-point detection (BOCPD) addresses non-stationarity by augmenting the state with the run length $r_t$, defined as the number of observations since the most recent change point \citep{Adams2007}. The learner maintains a separate posterior over $\theta$ conditional on each possible run length, which in the Bernoulli case is again a Beta density,
\begin{equation}
    p(\theta \mid r_t=r, x_{1:t})
    = \frac{1}{B(a_t^{(r)},b_t^{(r)})}\,\theta^{a_t^{(r)}-1}(1-\theta)^{b_t^{(r)}-1}.
\end{equation}
Here $a_t^{(r)}$ and $b_t^{(r)}$ are the success and failure counts the learner would have accumulated if the current run had length $r$, plus the prior pseudo-counts $a_0$ and $b_0$. They form the sufficient statistics under hypothesis $r$ and are much more compact than the full history.

The predictive probability of a success under a single run-length hypothesis is
\begin{equation}
    p(x_{t+1}=1 \mid r_t=r, x_{1:t})
    = \int_0^1 \theta\,p(\theta \mid r_t=r,x_{1:t})\,d\theta
    = \frac{a_t^{(r)}}{a_t^{(r)}+b_t^{(r)}},
\end{equation}
and the full Bayesian prediction averages these estimates over the run-length posterior,
\begin{equation}
    p(x_{t+1}=1 \mid x_{1:t})
    = \sum_r \frac{a_t^{(r)}}{a_t^{(r)}+b_t^{(r)}}\,p(r_t=r \mid x_{1:t}).
\end{equation}

The run-length posterior itself follows a simple recursive update. Between two observations the current run either continues, in which case $r_t = r_{t-1}+1$ and the Beta counts are updated with $x_t$, or a change point occurs, in which case $r_t = 0$ and the sufficient statistics reset to the prior pseudo-counts. The prior probability of a change is set by the hazard function
\begin{equation}
    H(r) = p(r_t=0 \mid r_{t-1}=r-1),
\end{equation}
which gives the probability that a run ends after surviving to age $r$. Throughout this chapter we use a constant hazard $H(r)=h$, corresponding to a geometric prior over regime durations; the probability of a change on the next trial is then independent of how long the current regime has already lasted.

Combining the predictive likelihood, the hazard and the sufficient statistics gives the joint update over $(r_t, a_t^{(r)}, b_t^{(r)})$ that is computed online. Figure~\ref{fig:cpd-bocpd-internals} shows the resulting internal variables for the example sequence of Fig.~\ref{fig:cpd-bernoulli-task}. The run-length posterior places its mass on a diagonal ridge that grows linearly during each stable regime, resets to small $r_t$ shortly after each change point, and occasionally bifurcates because the learner is initially uncertain which of the surviving run-length hypotheses is correct (Fig.~\ref{fig:cpd-bocpd-internals}A). The run-length-conditional means $a_t^{(r)}/(a_t^{(r)}+b_t^{(r)})$ track the latent $\theta_t$ at very different speeds: short run lengths follow recent observations closely and fluctuate, while long run lengths integrate over many trials and respond slowly (Fig.~\ref{fig:cpd-bocpd-internals}B). The marginal predictive after summing over run lengths recovers fast adaptation after each change point (Fig.~\ref{fig:cpd-bocpd-internals}C). Under constant hazard, the posterior mass on $r_t=0$ is uninformative because it equals the hazard at every trial; we therefore use the recent-run probability $p(r_t \leq 5 \mid x_{1:t})$ as a data-driven change indicator, which peaks shortly after each true change point (Fig.~\ref{fig:cpd-bocpd-internals}D).

\begin{figure}[htbp]
\centering
\includegraphics[width=0.92\linewidth]{cpd_bocpd_internals.pdf}
\caption{Bayesian online change-point detection for Bernoulli observations. The model maintains a posterior over run length and a Beta posterior for each run-length hypothesis. Panel A: heat map of $\log_{10} p(r_t\mid x_{1:t})$. Panel B: run-length conditional means $a_t^{(r)}/(a_t^{(r)}+b_t^{(r)})$ for selected $r$. Panel C: marginal predictive after summing over run lengths. Panel D: data-driven change-point indicator $p(r_t\leq 5\mid x_{1:t})$, which spikes after inferred change points (under constant hazard $p(r_t=0)$ is equal to the hazard itself, so the recent-run indicator is more informative).}
\label{fig:cpd-bocpd-internals}
\end{figure}

Although we use Bernoulli observations for clarity, the same logic generalizes to any conjugate exponential-family observation model: the counts $a_t^{(r)}$ and $b_t^{(r)}$ are simply the sufficient statistics of the Bernoulli likelihood and are replaced by the corresponding sufficient statistics for Gaussian, Poisson or other observation models.

\section{Approximation by a bank of fixed learning rates}
\label{sec:cpd-finite-bank}

The exact BOCPD posterior is supported on every possible run length up to the current trial. Any biologically plausible implementation is unlikely to maintain so many run-length-conditional estimates explicitly. A practical approach is to replace this dense posterior with a small bank of learners that each use a fixed learning rate, and to combine them through a downstream readout \citep{Nassar2010, Wilson2013}. We use this strategy here both as an approximation of BOCPD and as the conceptual bridge to the architectural question we ask in the next section: what does it take to build such a multi-timescale predictor in the brain?

We approximate the BOCPD posterior by a bank of $K$ fixed-rate learners
\begin{equation}
    \hat\theta_t^{(k)}, \qquad k=1,\ldots,K,
\end{equation}
each updated by a delta rule with its own fixed learning rate $\alpha_k$,
\begin{equation}
    \hat\theta_t^{(k)} = \hat\theta_{t-1}^{(k)} + \alpha_k\,(x_t - \hat\theta_{t-1}^{(k)}).
\end{equation}
Each $\alpha_k$ defines an effective integration timescale $\tau_k = 1/\alpha_k$ over which past observations contribute to that learner's estimate. A downstream readout combines the bank into a single prediction,
\begin{equation}
    \hat\theta_t = \sum_{k=1}^{K} w_t^{(k)}\,\hat\theta_t^{(k)},
    \qquad
    \sum_k w_t^{(k)} = 1.
\end{equation}
The readout weights $w_t^{(k)}$ are updated online by Bayesian filtering over which learner currently provides the best prediction (Sec.~\ref{sec:cpd-methods}).

We compare three predictors against exact BOCPD on the example sequence (Fig.~\ref{fig:cpd-finite-bank}): a single delta-rule learner with intermediate rate $\alpha = 0.1$; a three-rate bank with $\alpha \in \{1/4, 1/14, 1/82\}$, chosen so that fast, intermediate and slow timescales span the run lengths seen in the data; and the reduced mixture-of-delta-rules of \citet{Wilson2013}, which uses 18 log-spaced run-length nodes between $r=1$ and $r=100$.

The bank of learners shows the timescale separation that motivates the construction (Fig.~\ref{fig:cpd-finite-bank}B). The fast learner tracks every fluctuation in $\theta_t$ but is noisy within a regime; the slow learner is smooth but lags strongly after change points; the intermediate learner lies between the two. The Bayesian readout exploits this complementarity (Fig.~\ref{fig:cpd-finite-bank}C). Immediately after a change point, when the slow learner has been driven far from the new contingency, the fast learner has the highest predictive likelihood and the readout weights shift toward it. Once the environment has been stable for many trials the slow learner becomes more accurate, the readout shifts back toward it, and the intermediate learner tends to dominate at intermediate times. The resulting mixture prediction closely follows exact BOCPD and matches it in mean predictive negative log-likelihood (Fig.~\ref{fig:cpd-finite-bank}D, E). The reduced model of \citet{Wilson2013} achieves the best score in this example, while the single-rate baseline is clearly worse than the multi-rate predictors. A small bank of fixed learning rates therefore captures the qualitative behaviour of the full run-length posterior.

\begin{figure}[htbp]
\centering
\includegraphics[width=0.92\linewidth]{cpd_finite_bank.pdf}
\caption{A finite bank of fixed learning rates as an approximation to BOCPD. Panel A: latent $\theta_t$, observations and exact BOCPD prediction. Panel B: estimates $\hat\theta_t^{(k)}$ for fast, intermediate and slow learners. Panel C: Bayesian readout weights $w_t^{(k)}$ over the bank, shifting toward fast learners after change points and toward slow learners during stable periods. Panel D: mixture prediction compared with exact BOCPD, the reduced mixture of \citet{Wilson2013} with 18 log-spaced nodes between $r=1$ and $r=100$, and a single fixed-rate baseline. Panel E: mean predictive negative log-likelihood per trial for each model.}
\label{fig:cpd-finite-bank}
\end{figure}

\section{Diverse learning rates as an approximation basis}
\label{sec:cpd-diverse-rates}

The run-length posterior of exact BOCPD already has a natural interpretation in terms of evidence windows. A short run length means that only recent observations should be trusted, a long run length that a longer history can be integrated, and the Bayesian prediction is a mixture over all such windows weighted by their posterior probability. Replacing the dense run-length posterior with a finite bank of fixed-rate delta rules makes this interpretation concrete: each rate corresponds to one effective window, and the readout weights play the role of the run-length posterior on a reduced grid (Fig.~\ref{fig:cpd-rate-diversity}A--D).

The only ingredient this approximation requires from the underlying biology is a \emph{diversity of effective learning rates}, plus a readout that can weight them. Where this diversity lives is a separate question. One natural locus is the consolidation gradient from hippocampus to cortex proposed by classical theories of systems memory consolidation \citep{McClelland1995,Roxin2013} and elaborated by the parallel-pathway theory introduced in the previous chapter (Ch.~\ref{ch:ppt}). On this mapping, fast learners correspond to hippocampal pathways and slow learners to cortical ones (Fig.~\ref{fig:cpd-rate-diversity}E). The chapter argues that the learning-rate diversity required by approximate change-point inference is already present in this architecture: the same diversity that produces graded semantic generalization can simultaneously serve as an approximation basis for adaptive prediction in changing environments. We return to alternative biological loci for the same rate diversity in the discussion.

\begin{figure}[htbp]
\centering
\includegraphics[width=0.92\linewidth]{cpd_pathway_interpretation.pdf}
\caption{Diverse learning rates as an approximation basis for change-point inference. Panel A: exact BOCPD as a mixture over run lengths $r_t$. Panel B: bank of fast, intermediate and slow learners with estimates $\hat\theta_t^{(k)}$. Panel C: readout weights $w_t^{(k)}$, which shift toward fast learners after change points and toward slow learners during stable periods. Panel D: mixture prediction compared with exact BOCPD. Panel E: one possible biological mapping, in which the rate diversity is realised by PPT pathways spanning the hippocampal-to-cortical gradient.}
\label{fig:cpd-rate-diversity}
\end{figure}

\section{Simulation methods}
\label{sec:cpd-methods}

All figures in this chapter were generated from the same synthetic Bernoulli sequence and the model implementations described below.\footnote{Code: \texttt{change\_point\_detection\_from\_distributed\_learning\_rates}, branch \texttt{thesis-cpd-figures}, \texttt{PrEMCE-behavior/simulations/thesis\_figures/}.}

\subsection{Bernoulli sequence}

We generated a piecewise-constant Bernoulli sequence of length $T=300$, divided into five regimes of $60$ trials with true success probabilities $\theta \in \{0.85, 0.20, 0.50, 0.90, 0.30\}$, giving change points at trials $60, 120, 180, 240$. Observations $x_t$ were sampled independently from $\mathrm{Bernoulli}(\theta_t)$ within each regime. The same sequence (seed $7$) is used in all figures.

\subsection{Bayesian online change-point detection}

We ran the recursion of Sec.~\ref{sec:cpd-bocpd} on a run-length grid $r \in \{0, 1, \ldots, T\}$, with the Beta prior $\mathrm{Beta}(a_0=1, b_0=1)$, constant hazard $H(r) = h = 0.01$, and initial condition $p(r_0=0) = 1$. The joint update is implemented as the standard message-passing recursion of \citet{Adams2007}: predictive likelihoods $p(x_t \mid r_{t-1}=r, x_{1:t-1}) = a_{t-1}^{(r)}/(a_{t-1}^{(r)}+b_{t-1}^{(r)})$ for $x_t=1$, weighted by the hazard mass into the change-point bucket $r_t=0$ and by the survival mass into the growth bucket $r_t = r_{t-1}+1$, followed by renormalization and a Beta-counts update. Mass that would leave the grid at $r=T$ is absorbed at the boundary.

\subsection{Bank of fixed learning rates}

The three-rate bank uses learning rates $\alpha \in \{1/4, 1/14, 1/82\}$, equivalent to integration timescales $\tau = 1/\alpha \in \{4, 14, 82\}$ trials. All learners are initialised at $\hat\theta_0^{(k)} = 0.5$, matching the prior mean of $\mathrm{Beta}(1,1)$. Each learner's estimate is updated by the delta rule defined in the main text. The readout weights $w_t^{(k)}$ implement a leaky Bayesian filter over which learner currently best matches the data,
\begin{equation}
    \tilde w_t^{(k)} \;\propto\; p\!\left(x_t \,\big|\, \hat\theta_{t-1}^{(k)}\right) \, w_{t-1}^{(k)},
    \qquad
    \sum_k \tilde w_t^{(k)} = 1,
    \label{eq:cpd-bayes-filter}
\end{equation}
\begin{equation}
    w_t^{(k)} = (1-\lambda)\,\tilde w_t^{(k)} + \lambda\, \tfrac{1}{K},
    \label{eq:cpd-leak}
\end{equation}
with predictive likelihood $p(x_t \mid \hat\theta_{t-1}^{(k)}) = \hat\theta_{t-1}^{(k)}$ if $x_t=1$ and $1-\hat\theta_{t-1}^{(k)}$ if $x_t=0$. The leak coefficient $\lambda = 0.02$ pulls the weights weakly toward the uniform distribution at every step. Without it the weights converge irreversibly onto the single rate that happens to be best in the first regime; the leak allows the readout to re-explore the bank after subsequent change points. Initial weights are uniform, $w_0^{(k)} = 1/K$. The single-rate baseline in Fig.~\ref{fig:cpd-finite-bank} corresponds to $K=1$ and $\alpha = 0.1$.

\subsection{Reduced mixture of delta rules}

We implemented the reduced model of \citet{Wilson2013} for Bernoulli observations on $K=18$ run-length nodes $r^{(k)}$, spaced log-uniformly between $r=1$ and $r=100$ and rounded to the nearest integer. Each node maintains a delta-rule estimate $\hat\theta_t^{(k)}$ with effective learning rate
\begin{equation}
    \alpha_k = \frac{1}{r^{(k)} + a_0 + b_0},
    \label{eq:wilson-alpha}
\end{equation}
which reproduces the posterior-mean update of a Beta with effective sample size $r^{(k)}$. The node weights are propagated through a transition matrix
\begin{equation}
    T_{ij} = (1-h)\,\delta\!\left[i = \mathrm{nn}(r^{(j)}+1)\right] + h\,\delta\!\left[i = i_\text{reset}\right],
    \label{eq:wilson-T}
\end{equation}
where $\mathrm{nn}(r)$ returns the index of the node closest to run length $r$, $i_\text{reset}$ is the index of the smallest-run-length node, and $\delta[\cdot]$ is the Iverson bracket. The first term encodes the growth transition (the current run is extended by one and mapped to the nearest available node), the second the change-point reset. At each trial the weights are propagated by $\tilde w_t = T\,w_{t-1}$ and then reweighted by the per-node predictive likelihood, exactly as in Eq.~\ref{eq:cpd-bayes-filter} but with $\tilde w_t$ replacing $w_{t-1}$ on the right. The initial weight is placed entirely on the smallest-run-length node, matching the BOCPD initial condition.

\subsection{Predictive loss}

The mean predictive negative log-likelihood reported in Fig.~\ref{fig:cpd-finite-bank}E is
\begin{equation}
    \mathrm{NLL} = -\frac{1}{T}\sum_{t=1}^{T} \log p_t,
    \qquad
    p_t = \begin{cases} \hat\theta_t & \text{if } x_t = 1, \\ 1-\hat\theta_t & \text{if } x_t = 0, \end{cases}
\end{equation}
clipped to $[10^{-9}, 1-10^{-9}]$ for numerical stability.

\section{Discussion}

Adaptive learning in a changing environment requires that the learner decides how much of the past remains relevant before it updates. A surprising observation can reflect either noise within the current regime or a true change in the latent state, and the two interpretations call for very different revisions of the predictive distribution. BOCPD captures this requirement in its purest form: it maintains a posterior over the run length and predicts by mixing run-length-conditional estimates. Reading the Bernoulli model in this way makes the computational requirement explicit. Each evidence window contributes an estimate, and an adaptive prediction is a weighted combination of these estimates.

A small bank of fixed-rate learners with a Bayesian readout gives this requirement a biologically interpretable form. The bank reproduces the prediction performance of exact BOCPD on the Bernoulli task and recovers the same qualitative dynamics as the reduced mixture of \citet{Wilson2013}. The construction only assumes a diversity of effective learning rates plus a downstream readout that can weight them; it makes no commitment about where in the brain the rate diversity is realised. This is the part of the argument that is robust to anatomical detail and that the chapter advances as its main claim.

Where the rate diversity lives is then a separate, biologically richer question. The most natural locus, and the one that motivates the chapter, is the consolidation gradient from hippocampus to cortex \citep{McClelland1995,Roxin2013}, with the PPT-style mapping shown in Fig.~\ref{fig:cpd-rate-diversity}E (Ch.~\ref{ch:ppt}). However, the same approximation basis can be realised in other ways. Different cortical or subcortical areas with intrinsically different plasticity rates would do; so would a single area in which different neurons, subpopulations or synapse types span a range of effective time constants, or a recurrent network whose dynamics support multiple intrinsic timescales. We return to these within-area implementations in the chapter on representational drift (Ch.~\ref{ch:drift}) and in the recurrent-network chapter (Ch.~\ref{ch:rnn}), where rate-diverse predictors arise from a different mechanism but play a functionally related role.\todo{verify chapter labels ch:drift, ch:rnn}

Whatever its anatomical substrate, the rate-diverse bank is unlikely to carry the full change-point computation on its own. Exact BOCPD also requires a posterior over run lengths and, if the hazard is uncertain, additional beliefs over the hazard parameters. A biologically realistic decomposition therefore distributes the ideal-observer ingredients across systems: the rate-diverse bank provides candidate evidence-window estimates, while frontal, cingulate, striatal and neuromodulatory systems contribute the signals required to weight or gate them. Several neural results are consistent with such a distributed implementation. Activity in anterior cingulate cortex tracks environmental volatility and the value of information in changing environments \citep{Behrens2007}; functionally dissociable signals related to changepoint probability, relative uncertainty and reward influence learning rate \citep{McGuire2014}; and dorsomedial frontal cortex carries value estimates simultaneously at multiple timescales \citep{Meder2017}. These observations support the idea that the brain represents ingredients of adaptive learning across multiple regions rather than computing the run-length posterior in any single area.

The hazard function is an assumption about the environment, not a property of the observer, and changing it changes the inference problem rather than the rate-diverse approximation basis. A constant hazard treats every trial as equally likely to terminate a regime; a decreasing hazard expresses the belief that long-lived regimes are more likely to persist; an increasing hazard expresses the belief that regimes have typical lifetimes. Any age-dependent hazard $H(r)$ can be plugged into the same BOCPD recursion, because $r_t$ is the relevant state. When the hazard is fixed but unknown the observer can infer it jointly with the run length, and when it varies over time a hierarchical model is needed \citep{Wilson2010}.

The relation between heavy-tailed memory and heavy-tailed hazards deserves a separate clarifying remark. A power-law-like hazard is a statement about environmental durations, whereas a power-law-like memory kernel produced by a broad distribution of fixed learning rates is a property of the observer. The two can coincide in particular regimes but they are different objects. The useful claim for the present chapter is therefore weaker and more robust: a broad, log-spaced set of integration timescales provides a flexible approximation basis when the relevant environmental timescale is uncertain, regardless of the precise hazard.

Finally, the formal derivation assumes that all rate-specific learners receive the same online observation stream, which is a known simplification when the rate diversity is realised by a consolidation hierarchy. Replay can reshape the effective evidence kernel of each learner by selectively resampling past observations, for example emphasising post-change observations after an inferred change point or preserving older evidence during stable periods. Integrating replay with the rate-mixture view of change-point inference is a natural extension that is left to future work.

\section{Chapter summary}

Adaptive prediction in changing environments requires the learner to combine estimates over multiple possible evidence windows, as made explicit by the run-length posterior of BOCPD. In the Bernoulli case, these estimates are Beta posteriors with run-length-specific counts $a_t^{(r)}$ and $b_t^{(r)}$. A small bank of fixed-rate delta-rule learners with a Bayesian readout reproduces the predictive performance of exact BOCPD and the qualitative dynamics of the reduced mixture of \citet{Wilson2013}, with readout weights that shift toward fast learners after change points and toward slow learners during stable periods. The only architectural ingredient this requires is a diversity of effective learning rates, which can be realised across pathways or areas, as in PPT-style systems consolidation, or within an area by neurons or networks with diverse intrinsic timescales. This identifies learning-rate diversity as a biologically plausible approximation basis for change-point inference, and gives that diversity an additional normative role beyond graded semantic generalization.
